\chapter{Conclusiones y trabajo futuro}
\label{chap:conclusiones}
\section{Conclusiones}
Se realizó una comparación reproducible entre SVM y Bosque Aleatorio sobre características manuales de octrees y una implementación nativa de Net5, en ModelNet40 y resoluciones equivalentes de $32^3$ y $64^3$. La evaluación conjunta permitió relacionar la exactitud con el costo de preparación, predicción, memoria y almacenamiento. El cumplimiento experimental de los objetivos se apoyó en artefactos y registros versionados, y sus conclusiones se restringen al protocolo ejecutado.

En relación con el primer objetivo, se implementó el procesamiento desde mallas hasta representaciones densas y octrees, con raíz en nivel cero, muestreo de $20\,000$ puntos y serialización dispersa. La construcción y validación de los octrees proporcionaron la representación necesaria para implementar el enfoque clásico de clasificación mediante la extracción posterior de características. La muestra controlada de diez objetos produjo veinte registros correctos al considerar ambas resoluciones, sin fallos de integridad y con equivalencia geométrica verificada. Esta comprobación no constituyó una auditoría exhaustiva de los $12\,311$ modelos. El caso \texttt{chair\_0001} permitió ilustrar las etapas sin sustituir la evaluación del conjunto.

Para el segundo objetivo, se definieron características geométricas y jerárquicas de dimensiones 17 y 18, se estableció su contrato y se entrenaron los clasificadores clásicos. SVM obtuvo $75{,}93\,\%$ y $77{,}19\,\%$ de exactitud; el bosque obtuvo $76{,}38\,\%$ y $77{,}19\,\%$. Dentro del protocolo común de los métodos clásicos, SVM presentó menor tamaño de archivo, latencia y memoria observada, mientras que el bosque no mostró una ventaja de exactitud que compensara sistemáticamente su mayor costo de almacenamiento.

Para el tercer objetivo, se configuró y entrenó Net5 sobre la representación jerárquica nativa. Sus exactitudes de $82{,}21\,\%$ en $32^3$ y $83{,}23\,\%$ en $64^3$ fueron superiores a las de los métodos clásicos. Esta ventaja también fue respaldada por los contrastes emparejados frente a SVM y Bosque Aleatorio en ambas resoluciones, con valores $p<10^{-11}$ (Tabla~\ref{tab:contrastes-finales}). Net5 alcanzó además los mayores F1 macro, aunque presentó confusiones recurrentes entre categorías y una diferencia entre las exactitudes de ajuste y validación. La similitud geométrica entre algunos pares de categorías constituye una explicación plausible de esas confusiones; no se demostró como su causa mediante un experimento controlado.

El cuarto objetivo se sustentó en la organización de configuraciones, particiones, versiones, modelos y registros por objeto. La reevaluación conservó las seis secuencias de predicciones, lo que respalda la reproducibilidad de la evaluación de los modelos almacenados. Esta evidencia no equivale a reproducir nuevos entrenamientos con distintas semillas ni a eliminar la variación temporal del entorno.

El quinto objetivo permitió identificar el compromiso entre desempeño y costo. Aumentar la resolución produjo mejoras de $0{,}81$ a $1{,}26$ puntos porcentuales de exactitud. Entre los contrastes por cambio de resolución, la diferencia fue nominalmente significativa únicamente para SVM, con $p=0{,}029$; no lo fue para Bosque Aleatorio, con $p=0{,}161$, ni para Net5, con $p=0{,}109$ (Tabla~\ref{tab:contrastes-finales}). No se aplicó corrección por comparaciones múltiples, por lo que la significación del SVM debe interpretarse con ese alcance. Los resultados no significativos no demuestran igualdad entre resoluciones.

La comparación de costos mostró que el aumento de resolución exigió un aumento del tiempo de entrenamiento, la latencia y el uso de memoria de Net5, aunque se conservaron el número de parámetros y el tamaño del modelo. En el enfoque clásico, SVM ofreció un compromiso favorable de latencia, memoria y almacenamiento frente al bosque bajo su protocolo común. Net5 en $32^3$ mantuvo una ventaja de exactitud sobre los clásicos en $64^3$ con menor costo de ejecución que su propia variante de mayor resolución. Estos resultados muestran que el tamaño del archivo no determina los recursos de ejecución y que una exactitud máxima no implica el mejor compromiso global. La elección debe considerar conjuntamente entrenamiento, inferencia, memoria pico y almacenamiento, respetando los alcances de medición: la extracción HCE fue estimada y los tiempos de ambos enfoques no abarcaron intervalos idénticos. Por ello, no se estableció una superioridad directa de velocidad entre HCE y Net5 ni se equipararon RAM, VRAM y tamaño en disco.

En el contexto del Grupo de Óptica y Tratamiento de Señales de la Escuela de Física de la Universidad Industrial de Santander, el trabajo aporta una base experimental para estudiar representaciones de información tridimensional. Su alcance comprobado fue la clasificación de objetos CAD; la transferencia a reconstrucción o metrología óptica requiere validación adicional con mediciones reales.

\section{Principales aportes}
Se integró un flujo trazable que relacionó mallas, representaciones, características, modelos y resultados. La definición de los descriptores HCE de 17 y 18 componentes, junto con sus contratos de dimensión, orden y significado, permitió documentar la información geométrica y jerárquica utilizada por los clasificadores clásicos. Se implementaron el enfoque clásico y Net5 con particiones y resoluciones equivalentes, lo que permitió comparar su desempeño de clasificación bajo una base común.

Se realizó una evaluación emparejada sobre los mismos $2\,468$ objetos oficiales de prueba y se conservaron las predicciones por identificador para analizar aciertos, errores y discordancias. Esta evaluación aportó evidencia cuantitativa del compromiso entre exactitud y costo: las mejoras de clasificación al aumentar la resolución fueron limitadas frente al incremento de varios indicadores de recursos. La separación entre tiempo de entrenamiento, tiempo de inferencia, memoria pico y tamaño del modelo permitió interpretar cada magnitud según su alcance, sin atribuir una única noción de eficiencia a medidas diferentes.

La conservación de modelos, particiones, predicciones, configuraciones, versiones, huellas y registros experimentales constituyó un aporte de trazabilidad y permitió reevaluar los modelos almacenados. Se documentaron criterios de equivalencia geométrica y persistencia, y se distinguieron los ejemplos ilustrativos, la muestra controlada y la evaluación completa de prueba. Esta organización permite relacionar las conclusiones con los artefactos que las respaldan y con sus limitaciones documentadas.

\section{Limitaciones}
El estudio se restringió a ModelNet40, objetos CAD, dos resoluciones, un equipo y una ejecución de entrenamiento por configuración. No se evaluaron otras semillas, otros conjuntos de datos, escaneos reales ni estudios de ablación; tampoco se midió el consumo energético. En consecuencia, no se estableció la estabilidad de los resultados entre entrenamientos ni su generalización a otros dominios de adquisición. La validación geométrica controlada comprendió únicamente diez objetos y no constituyó una auditoría exhaustiva de los $12\,311$ modelos.

La muestra de memoria abarcó tres categorías y los tiempos GPU y CPU tuvieron una diferencia de instrumentación en el armado del lote. La generación inicial de octrees no formó parte del tiempo de inferencia y el costo de extracción HCE de entrenamiento fue estimado. Estas diferencias impiden interpretar todos los tiempos como mediciones del mismo intervalo. Los registros complementarios fila--objeto de HCE coincidieron con los identificadores y posiciones del manifiesto de Net5 en ambas resoluciones. La coincidencia única de descriptores se documentó mediante los
registros suministrados y la revisión del código de verificación
posterior, sin repetir de forma independiente su ejecución sobre
los octrees originales. Los contrastes estadísticos se interpretaron sin ajuste por multiplicidad. Estas condiciones delimitan la evidencia y no se sustituyeron por afirmaciones de generalidad o de equivalencia entre métodos.

\section{Trabajo futuro}
Para abordar la limitación de un entrenamiento por configuración, se propone repetir los entrenamientos con varias semillas, cuantificar la variabilidad entre ejecuciones y diseñar contrastes con control de multiplicidad. El estudio de más resoluciones permitiría examinar si las ganancias y los costos observados se mantienen fuera del intervalo evaluado. La incorporación de otros conjuntos de datos y equipos permitiría valorar la generalización de la clasificación y la dependencia del rendimiento respecto del entorno computacional.

La validación geométrica puede ampliarse desde la muestra de diez objetos hasta el conjunto completo, conservando los registros de equivalencia e integridad por objeto. La medición de recursos puede extenderse a muestras equilibradas de las cuarenta categorías y a intervalos simétricos desde la malla o desde el NPZ, incorporando el armado de lotes en todas las variantes y cronometrando la extracción HCE completa. También conviene separar lecturas con y sin caché, evaluar otros tamaños de lote y números de hilos, e incorporar mediciones energéticas para estudiar una dimensión de costo no cubierta en este trabajo.

La ausencia de estudios de ablación motiva evaluar de forma controlada la contribución de los grupos de descriptores HCE, los atributos de entrada y los componentes de Net5. Estos estudios permitirían contrastar las explicaciones propuestas sobre las diferencias de desempeño. La preparación geométrica de Net5 constituye además una etapa concreta para optimización mediante reutilización de planes y reducción de transferencias, cuya mejora debe verificarse sin alterar las predicciones. Asimismo, pueden estudiarse otras capacidades de red, descriptores complementarios y estrategias de combinación de clasificadores.

Finalmente, para superar la restricción a objetos CAD, se propone una evaluación posterior sobre reconstrucciones o mediciones tridimensionales reales, incluidas las obtenidas mediante técnicas ópticas. Esta extensión requiere datos con ruido, oclusiones y condiciones de adquisición documentadas. Solo mediante esa validación adicional podrá establecerse la utilidad de las representaciones estudiadas dentro de una tarea específica de metrología óptica.
